\documentclass[11pt]{article}

\usepackage[margin=1in]{geometry}
\usepackage{amsmath,amssymb,amsthm}
\usepackage{booktabs}
\usepackage{natbib}
\usepackage{hyperref}
\hypersetup{colorlinks=true,linkcolor=blue,citecolor=blue,urlcolor=blue}

\newtheorem{proposition}{Proposition}
\newtheorem{lemma}{Lemma}
\newtheorem{remark}{Remark}

\newcommand{\eps}{\varepsilon}
\newcommand{\Var}{\operatorname{Var}}
\newcommand{\Cov}{\operatorname{Cov}}
\newcommand{\sinck}{\operatorname{\mathsf{sinc}}}
\newcommand{\coef}[1]{[u^{#1}]}

\title{A Poissonized Evaluator for the Multinomial Inversion Series\\
\large Exact de-Poissonization by a discrete Cauchy integral on the saddlepoint circle}
\author{}
\date{}

\begin{document}
\maketitle

\begin{abstract}
Every term of the inversion series is a constrained transform: the coefficient of $u^N$ in a
product of $K$ per-category polynomials. The existing evaluators obtain it by multiplying the
polynomials together and cutting the running product back to degree $N$ after every category; in
the FFT tree that truncation fixes the working size at $M\approx2$--$4N$ and costs a full descent,
leaf convolution and ascent per category, driven by a precomputed plan of merge factors whose size
grows like $KM\log M$. This document describes an evaluator that drops the constraint instead of
enforcing it. The categories become independent, each polynomial is evaluated on its own at $Q$
equispaced points of a circle, the $K$ evaluations are multiplied pointwise, and the coefficient is
recovered \emph{exactly}, up to aliasing, by a $Q$-point trapezoid rule for Cauchy's integral. With
the circle's radius at the saddlepoint of the tilted total count --- the same tilt the gamma search
and the FFT tree already use --- the aliasing is suppressed like the tails of that count, and two
facts make the choice of $Q$ certified rather than heuristic: the aliasing at frequency zero bounds
it for every term of the series, and that frequency-zero aliasing can be measured exactly by
comparing $Q$ against $2Q$. Per frequency the work is one fold of $N+1$ coefficients and one
length-$Q$ FFT per category, with $Q$ of seven to ten standard deviations of the tilted count. It
agrees with the FFT tree to $10^{-15}$--$2.3\times10^{-12}$ on the p-value at a fixed term budget.
Measured on an Apple M5, one term costs about a third as much in the deep tail when $Q$ equals the
tree's working size ($2.1$--$3.6\times$ faster), $4$--$20\times$ less when $Q$ is a half or a
quarter of it, which happens at many categories, and $4$--$22\times$ less away from the tail; the
plan is $30$--$140\times$ smaller ($4$\,MB against $588$\,MB at $N=10^4$, $K=50$).

The same idea -- drop the constraint, work with independent categories -- is already used
elsewhere in this codebase to select the tilt $\gamma$ itself, and we derive that construction in
full: the Poissonized log-MGF is convex in $\gamma$, its derivatives come out of the same
one-pass moment computation the tilt $\nu$ is already solved for by an application of the envelope
theorem, and a safeguarded Newton iteration replaces a 128-point grid at $13$--$25\times$ less time
per instance. We explain why the same \emph{approximate} (never de-Poissonized) construction that
is exact enough for $\gamma$ is not accurate enough for the series terms themselves. Three further,
evaluator-independent changes round out the implementation: a persistent thread pool replaces
per-chunk thread creation in \texttt{series()} ($2\times$ at the smallest instances); the adaptive
stopping rule for \texttt{extrapolate}$\,\ge2$ combines the pre-existing term-ratio rule with a
check on a shallow, median-smoothed extrapolation's trailing-window stability, verified stable to
$10^{-9}$--$10^{-13}$ under a one-ulp perturbation of every term, across both evaluators and all
thread counts tried; and locating the statistic's range before the series can run -- previously an
$O(KN^2)$ dynamic program -- is now $O(K)$ for the power-divergence family (an exact vertex
argument for the minimum, a closed-form bound for the maximum) and $O(K\log K)$ for the
probability-mass statistic (an exact maximum via \citet{whitehendy2010}'s multinomial-mode
algorithm), cutting setup time by up to $4$--$5\times$ at $N=10^4$.
\end{abstract}

\tableofcontents

\section{What each series term is}
\label{sec:setup}

We keep the notation of the companion documents. $X\sim\mathrm{Multinomial}(N,\mathbf p)$ over $K$
categories, $S=\sum_k a_k(X_k)$ is the centred statistic, $\gamma$ is the tilt chosen by
\texttt{fill\_gamma}, $z_n=\gamma+inT$, and $M_S(z)=\mathbb E[e^{-zS}]$. The series accumulates
$t_n=M_S(z_n)\sinck(z_n)$ with $\sinck(z)=\mathbb E[e^{-zU}]$, $U\sim\mathrm{Unif}[0,W]$, $W=\pi/T$;
see \texttt{error\_bounds.tex} for why $p=\tfrac12\sum_{n\in\mathbb Z}t_n$. Everything below concerns
computing $M_S(z_n)$ for one frequency $n$.

Write $p_k(i)=p_k^i/i!$, so that $\log p_k(i)=i\log p_k-\log i!$ as in \texttt{rates.tex}. After
\texttt{fill\_cells}, cell \texttt{matrix[k*(N+1)+i]} stores
\[
  \ell_k(i)\;:=\;\texttt{log\_prob}\;=\;\log p_k(i)-\gamma a_k(i),
  \qquad
  \theta_k(i)\;:=\;\texttt{reward}\;=\;T\,a_k(i),
\]
so that $p_k(i)\,e^{-z_na_k(i)}=e^{\ell_k(i)}e^{-in\theta_k(i)}$. Define the per-category polynomials
\[
  \phi_{k,n}(u)\;=\;\sum_{i=0}^{N}e^{\ell_k(i)}\,e^{-in\theta_k(i)}\,u^i .
\]

\begin{proposition}[The term as a coefficient]
\label{prop:coef}
For every $n$,
\[
  M_S(z_n)\;=\;N!\;\coef{N}\prod_{k=1}^{K}\phi_{k,n}(u).
\]
\end{proposition}
\begin{proof}
$P(X=x)=N!\prod_k p_k(x_k)$ on $\{x\in\mathbb Z_{\ge0}^K:\sum_kx_k=N\}$, so
$M_S(z_n)=N!\sum_{\sum x_k=N}\prod_k e^{\ell_k(x_k)-in\theta_k(x_k)}$. Expanding the product of the
$\phi_{k,n}$, the coefficient of $u^N$ collects exactly these products; the bound $x_k\le N$ is
automatic.
\end{proof}

\section{The truncated product, and what it costs}
\label{sec:tree}

The reference evaluator (\texttt{eval\_fourier\_thread}) and the FFT tree
(\texttt{fourier\_fft\_precompute.h}, analysed in \texttt{rates.tex}) both build the product one
category at a time and discard every coefficient above degree $N$ after each step --- those never
reach $\coef N$. Truncation keeps every intermediate polynomial at degree $\le N$, so each step is a
product of two polynomials of degree $\le N$ and fits in a ring of size
$M=2\cdot2^{\lceil\log_2(N+1)\rceil}$.

That truncation is also what makes each step expensive. The tree works modulo $x^M-1$ by
recursive ring splitting: a descent to width $4$, direct convolutions at the leaves, and an ascent
that reconstructs the product before it can be truncated. Counting call sites in
\texttt{fourier\_fft\_count\_calls},
\begin{equation}
  \label{eq:calls}
  \mathrm{calls}(M)\;=\;\tfrac32M\,(\log_2M-2)\;+\;4M
\end{equation}
merges per category, each applying a stored magnitude factor and a complex rotation to all SIMD
lanes; the precompute plan stores $K\cdot\mathrm{calls}(M)$ factors of 16 bytes. At $N=1000$ that is
$M=2048$ and $35{,}840$ merges per category; at $N=10^4$ it is $M=32{,}768$, $770{,}048$ merges, and
a $588$\,MB plan for $K=50$.

\section{Poissonization}
\label{sec:poisson}

The constraint $\sum_kx_k=N$ is what couples the categories. The classical way to remove it is to
make the counts independent Poisson variables and condition on their total afterwards: if
$Y_k\sim\mathrm{Poisson}(\lambda p_k)$ independently, then $(Y_1,\dots,Y_K)\mid\sum_kY_k=N$ is
$\mathrm{Multinomial}(N,\mathbf p)$ \citep{levin1981}. In generating-function language, the product
$\prod_k\phi_{k,n}$ is the unconstrained (Poissonized) object, and $\coef N$ is the conditioning; the
whole difficulty lies in undoing the Poissonization, the \emph{de-Poissonization}
\citep{jacquetszpankowski1998}.

\paragraph{The tilt.} For $\nu\in\mathbb R$ and $\rho=e^\nu$, define on $\{0,\dots,N\}$
\[
  w_k(i)=e^{\ell_k(i)+\nu i},\qquad Z_k=\sum_{i=0}^Nw_k(i),\qquad \pi_k(i)=\frac{w_k(i)}{Z_k},
\]
and let $Y_k\sim\pi_k$ be independent with total $T=\sum_kY_k$. Then
\begin{equation}
  \label{eq:pgf}
  \prod_k\phi_{k,0}(\rho u)\;=\;\Big(\prod_kZ_k\Big)\,\mathbb E\big[u^T\big],
  \qquad\text{so}\qquad
  \coef N\prod_k\phi_{k,0}(u)\;=\;\rho^{-N}\Big(\prod_kZ_k\Big)P(T=N).
\end{equation}
We fix $\nu$ by $\mathbb E[T]=N$. The solution is unique, since
$\frac{d}{d\nu}\mathbb E[T]=\Var T>0$, and it is precisely the saddlepoint: it minimizes
$\rho^{-N}\prod_kZ_k$ over $\rho$. It is the same tilt that \texttt{rates.tex} derives for the FFT
tree (\texttt{l\_rate}) and that the Newton gamma search solves for in \texttt{gamma\_search.h}; the
evaluator re-solves it on the transformed cells in \texttt{poisson\_tilt}, warm-started from
\texttt{mult->l\_rate}, so it does not depend on which search produced $\gamma$. Throughout,
\[
  V\;:=\;\Var T\;=\;\sum_k\Var Y_k
\]
is the variance of the tilted total count, the \texttt{v\_i} accumulator of \texttt{gamma\_moments}.

\section{Exact de-Poissonization: the discrete Cauchy formula}
\label{sec:cauchy}

\begin{proposition}[Trapezoid rule for a Taylor coefficient]
\label{prop:cauchy}
Let $F(u)=\sum_{j=0}^{D}c_ju^j$, $\rho>0$, $Q\ge1$ and $\omega=e^{2\pi i/Q}$. Then
\[
  \frac1Q\sum_{q=0}^{Q-1}\omega^{-qN}\rho^{-N}F(\rho\,\omega^q)
  \;=\;\sum_{\ell\in\mathbb Z}c_{N+\ell Q}\,\rho^{\ell Q},
\]
with $c_j:=0$ outside $[0,D]$. In particular the left side equals $c_N$ exactly whenever
$Q>\max(N,D-N)$.
\end{proposition}
\begin{proof}
Substituting the expansion, $\frac1Q\sum_q\omega^{q(j-N)}$ is $1$ if $j\equiv N\pmod Q$ and $0$
otherwise. A term with $\ell\ge1$ needs $N+\ell Q\le D$, i.e.\ $Q\le D-N$; one with $\ell\le-1$ needs
$N+\ell Q\ge0$, i.e.\ $Q\le N$.
\end{proof}

This is the discretized Cauchy integral over the circle $|u|=\rho$, the standard device for
computing Taylor coefficients numerically \citep{lynessmoler1967,bornemann2011}; for analytic
periodic integrands the trapezoid rule converges exponentially \citep{trefethenweideman2014}, and
here the integrand is a trigonometric polynomial, so the only error is the aliasing written out
explicitly above.

Applied to $F=\prod_k\phi_{k,n}$, which has degree $D=KN$, Propositions~\ref{prop:coef}
and~\ref{prop:cauchy} give
\begin{equation}
  \label{eq:main}
  M_S(z_n)+A_n(Q)\;=\;
  \underbrace{N!\,\rho^{-N}\prod_kZ_k}_{\texttt{resN\_exp}}\cdot
  \underbrace{\frac1Q\sum_{q=0}^{Q-1}\omega^{-qN}\prod_{k=1}^{K}\hat\phi_{k,n}(q)}_{\texttt{unit}},
  \qquad
  \hat\phi_{k,n}(q)=\sum_{i=0}^{N}\pi_k(i)\,e^{-in\theta_k(i)}\,\omega^{qi},
\end{equation}
where the aliasing is
\[
  A_n(Q)\;=\;N!\,\rho^{-N}\sum_{\ell\ne0}c_{N+\ell Q}(n)\,\rho^{N+\ell Q},
  \qquad c_j(n):=\coef{j}\prod_k\phi_{k,n}(u),
\]
and vanishes identically once $Q>\max(N,(K-1)N)$.

Three features of~\eqref{eq:main} are what make it cheap.
\begin{itemize}
  \item \textbf{No coupling.} Each $\hat\phi_{k,n}$ is computed from category $k$ alone, and the
        categories meet only in a pointwise product over $q$.
  \item \textbf{One FFT per category.} Folding the coefficients modulo $Q$,
        $\hat\phi_{k,n}(q)=\sum_{b<Q}B_{k,n}(b)\,\omega^{qb}$ with
        $B_{k,n}(b)=\sum_{i\equiv b\,(Q)}\pi_k(i)e^{-in\theta_k(i)}$, so all $Q$ values come from a single
        length-$Q$ transform of the folded vector.
  \item \textbf{Modulus at most one.} Since $\sum_i\pi_k(i)=1$, every $|\hat\phi_{k,n}(q)|\le1$ and so
        is every partial product. Everything that could overflow or underflow --- $N!$, $\rho^{-N}$
        and the $Z_k$ --- sits in the frequency-independent factor, which the plan computes once in
        quad precision, exactly as the precompute plan does for its own \texttt{resN\_exp}.
\end{itemize}

\begin{remark}[Conditioning of the extraction]
At $n=0$ the terms of the $q$-sum have modulus at most one while their mean is $P(T=N)$, which by
the local central limit theorem is about $1/\sqrt{2\pi V}$. The extraction therefore loses about
$\log_{10}\sqrt{2\pi V}$ digits to cancellation: $3.0$ digits at $N=1000$, $K=7$, and $3.8$ digits
at $N=10^4$, $K=3$. This is the sense in which the saddlepoint radius is the natural one: it
minimizes $\rho^{-N}\prod_kZ_k$, the modulus bound the cancellation has to reduce
\citep{bornemann2011}.
\end{remark}

\section{Aliasing: one bound for every term, and a certificate for it}
\label{sec:alias}

\begin{lemma}[Frequency zero dominates]
\label{lem:dominate}
For all $n$ and $j$, $|c_j(n)|\le c_j(0)$. Consequently, with
\[
  \alpha(Q)\;:=\;\sum_{\ell\ne0}\frac{P(T=N+\ell Q)}{P(T=N)},
\]
the aliasing satisfies $|A_n(Q)|\le A_0(Q)=\alpha(Q)\,M_S(\gamma)$, and the aliasing in the series
terms satisfies $|\Delta t_n|\le\alpha(Q)\,t_0$ for every $n$.
\end{lemma}
\begin{proof}
$c_j(n)=\sum_{i_1+\dots+i_K=j}\prod_ke^{\ell_k(i_k)}e^{-in\theta_k(i_k)}$, and $c_j(0)$ is the same sum
of moduli, which gives the first claim by the triangle inequality. By~\eqref{eq:pgf},
$c_j(0)\rho^j=\prod_kZ_k\,P(T=j)$, so $A_0(Q)=N!\rho^{-N}\prod_kZ_k\sum_{\ell\ne0}P(T=N+\ell Q)$ while
$M_S(\gamma)=N!\rho^{-N}\prod_kZ_k\,P(T=N)$. Finally
$|\sinck(z_n)|=|\mathbb E\,e^{-z_nU}|\le\mathbb E\,e^{-\gamma U}=\sinck(\gamma)$.
\end{proof}

So a single number, $\alpha(Q)$, bounds the aliasing of every term on the absolute scale of the first
one. The bound is not merely valid but attained: in every case we measured, the largest
term-by-term aliasing over the first 256 frequencies equals $\alpha(Q)$ to two or three significant
digits (Table~\ref{tab:cert}), i.e.\ the worst term is $t_0$ itself.

\paragraph{The certificate.} $\alpha(Q)$ is measurable without knowing $P(T=N)$. At $n=0$ the unit in
\eqref{eq:main} is $\sigma(Q):=\sum_{\ell}P(T=N+\ell Q)$, since
$\prod_k\hat\phi_{k,0}(q)=\mathbb E[\omega^{qT}]$. Comparing $Q$ with $2Q$ removes the even multiples:

\begin{proposition}[What comparing $Q$ with $2Q$ measures]
\label{prop:cert}
With $D(Q):=\big(\sigma(Q)-\sigma(2Q)\big)/\sigma(2Q)$,
\[
  D(Q)\;=\;\frac{\alpha(Q)-\alpha(2Q)}{1+\alpha(2Q)} .
\]
Hence $D(Q)\le\alpha(Q)$ always. If moreover $\alpha(2Q)\le\alpha(Q)^2$, then
$\alpha(Q)\le D(Q)\,\frac{1+\alpha(Q)^2}{1-\alpha(Q)}=D(Q)\,\big(1+O(D(Q))\big)$.
\end{proposition}
\begin{proof}
$\sigma(Q)-\sigma(2Q)=\sum_{\ell\text{ odd}}P(T=N+\ell Q)=P(T=N)\,\big(\alpha(Q)-\alpha(2Q)\big)$ and
$\sigma(2Q)=P(T=N)\,(1+\alpha(2Q))$. For the second part,
$\alpha(Q)-\alpha(Q)^2\le\alpha(Q)-\alpha(2Q)=D(Q)(1+\alpha(2Q))\le D(Q)(1+\alpha(Q)^2)$.
\end{proof}

The condition $\alpha(2Q)\le\alpha(Q)^2$ says that doubling $Q$ at least squares the tail
suppression. It holds whenever the law of $T$ is log-concave: then $d\mapsto\log P(T=N+d)/P(T=N)$ is
concave and vanishes at $d=0$, so $P(T=N+2\ell Q)/P(T=N)\le\big(P(T=N+\ell Q)/P(T=N)\big)^2$ for each
$\ell$, and summing over $\ell\ne0$ gives the claim.

\begin{proposition}[Log-concavity for the probability-mass statistic]
\label{prop:logconcave}
For the probability-mass statistic with $\gamma\le1$, every $\pi_k$ is log-concave on
$\{0,\dots,N\}$, and so is the law of $T$.
\end{proposition}
\begin{proof}
Here $a_k(i)=\log p_k(i)+\mathrm{const}_k$, so $\ell_k(i)+\nu i=(1-\gamma)(i\log p_k-\log i!)+\nu i+\mathrm{const}$.
Its second difference in $i$ is $-(1-\gamma)\log\frac{i+1}{i}\le0$, because linear terms have none
and $\Delta^2(-\log i!)=-\log\frac{i+1}{i}$. The support is an interval, so $\pi_k$ is log-concave,
and convolutions of log-concave sequences without internal zeros are log-concave
\citep{hoggar1974}.
\end{proof}

\begin{remark}[The log-likelihood ratio]
For LLR, $a_k(i)=-i\log\frac{i}{Np_k}+\mathrm{const}_k$ and
$\ell_k(i)+\nu i=-\log i!+\gamma\,i\log i+(\text{affine in }i)$. Since
$\Delta^2(-\log i!)=-\log\frac{i+1}{i}=-\frac1i+\frac1{2i^2}+O(i^{-3})$ and
$\Delta^2(i\log i)=\frac1i+O(i^{-3})$, its second difference is
$-(1-\gamma)/i+1/(2i^2)+O(i^{-3})$. It is log-concave only for
$i\gtrsim1/(2(1-\gamma))$ --- about $i>125$ at $\gamma=0.996$ --- so
Proposition~\ref{prop:logconcave} does not cover the deep tail. For LLR the certificate rests on
measurement instead, and Table~\ref{tab:cert} shows it matching the actual aliasing there just as
closely.
\end{remark}

\paragraph{Choosing $Q$.} \texttt{build\_poisson\_plan} proceeds as follows.
\begin{enumerate}
  \item Start at the smallest power of two $Q\ge4\sqrt V$ (and $\ge16$). If
        $Q>\max(N,(K-1)N)$ the formula is exact and we stop.
  \item Compute $D(Q)$ from $\sigma(Q)$ and $\sigma(2Q)$, in a separate scalar code path
        (\texttt{poisson\_s0}) that shares nothing with the vectorized evaluator.
  \item Accept $Q$ if $D(Q)\le10^{-10}$, or if $D(Q)<10^{-6}$ has stopped decreasing by at least a
        factor ten --- the rounding floor, where aliasing is already gone. The smallness condition
        matters: below about two standard deviations $D$ saturates near one and also stops
        decreasing.
  \item Otherwise double $Q$ and repeat, stopping as soon as $Q>\max(N,(K-1)N)$; $Q$ is capped at
        $2^{20}$ with a warning.
\end{enumerate}
Starting low matters because $Q$ moves in factors of two and the certificate can only raise it. The
tolerance $10^{-10}$ applies to the \emph{terms}, which is conservative for $p$: aliasing
concentrates in the lowest frequencies and largely cancels in the sum, and its measured effect on
the partial sums is two to six orders of magnitude below $D(Q)$ (last column of
Table~\ref{tab:cert}).

\begin{table}[t]
\centering
\small
\begin{tabular}{llrrccc}
\toprule
instance & $\sqrt V$ & $Q$ & $Q/\sqrt V$ & certificate $D(Q)$ & $\max_{n<256}|\Delta t_n|/t_0$ & partial-sum error \\
\midrule
PMF $N=1000$, $K=5$  & 237   & 512   & 2.2  & $1.85\times10^{-1}$ & $1.86\times10^{-1}$ & $7.1\times10^{-6}$ \\
                     &       & 1024  & 4.3  & $6.96\times10^{-4}$ & $6.96\times10^{-4}$ & $1.1\times10^{-9}$ \\
                     &       & 2048  & 8.6  & $3.80\times10^{-11}$ & $3.80\times10^{-11}$ & $2.3\times10^{-13}$ \\
                     &       & 4096  & 17.3 & $2.76\times10^{-14}$ & $2.76\times10^{-14}$ & $6.4\times10^{-16}$ \\
\midrule
LLR $N=1000$, $K=7$, $p_k\propto k^2$ & 399 & 1024 & 2.6 & $6.35\times10^{-2}$ & $6.36\times10^{-2}$ & $3.5\times10^{-7}$ \\
                     &       & 2048  & 5.1  & $1.63\times10^{-4}$ & $1.63\times10^{-4}$ & $1.2\times10^{-10}$ \\
                     &       & 4096  & 10.3 & $9.71\times10^{-14}$ & $9.70\times10^{-14}$ & $6.2\times10^{-14}$ \\
\midrule
LLR $N=10^4$, $K=10$ & 1989  & 4096  & 2.1  & $2.36\times10^{-1}$ & $2.37\times10^{-1}$ & $1.2\times10^{-4}$ \\
                     &       & 8192  & 4.1  & $8.96\times10^{-4}$ & $8.96\times10^{-4}$ & $3.7\times10^{-10}$ \\
                     &       & 16384 & 8.2  & $7.59\times10^{-11}$ & $7.60\times10^{-11}$ & $2.5\times10^{-14}$ \\
                     &       & 32768 & 16.5 & $3.82\times10^{-13}$ & $4.31\times10^{-13}$ & $2.9\times10^{-14}$ \\
\midrule
PMF $N=10^4$, $K=5$  & 2358  & 8192  & 3.5  & $6.78\times10^{-3}$ & $6.78\times10^{-3}$ & $7.0\times10^{-7}$ \\
                     &       & 16384 & 6.9  & $7.16\times10^{-8}$ & $7.16\times10^{-8}$ & $2.8\times10^{-11}$ \\
                     &       & 32768 & 13.9 & $4.43\times10^{-13}$ & $4.43\times10^{-13}$ & $1.3\times10^{-15}$ \\
\bottomrule
\end{tabular}
\caption{The certificate against the aliasing actually present. $Q$ is forced; the reference is the
exact evaluation with $Q>\max(N,(K-1)N)$, over the first 256 frequencies. Uniform
$x_0=N/K$ and $p_k\propto k$ unless stated. The certificate reproduces the worst-case term aliasing to two or
three digits (Lemma~\ref{lem:dominate} is attained at $n=0$), and the partial sum
$\tfrac12t_0+\sum_{0<n<256}t_n$ is disturbed far less. Rows selected by
\texttt{build\_poisson\_plan} under the $10^{-10}$ tolerance: $2048$, $4096$, $16384$ and $32768$
respectively.}
\label{tab:cert}
\end{table}

\section{Why $Q$ is large exactly in the deep tail}
\label{sec:variance}

Since $Q$ scales like $\sqrt V$ while $M$ scales like $N$, everything turns on how broad the tilted
per-category laws are.

\paragraph{Away from the tail} $\gamma$ is small, the tilt is mild, and $\pi_k$ is close to a
truncated Poisson law: at $\gamma=0$, $\pi_k(i)\propto(p_k\rho)^i/i!$, the constraint
$\mathbb E[T]=N$ gives $\rho=N$, and $V=N$ exactly. Measured with $x_0$ at the mode: $\sqrt V=10.8$
at $N=100$, $K=20$ ($\gamma=0.14$), and $\sqrt V=33$--$37$ at $N=1000$. There $Q\approx10\sqrt N\ll M$,
and the fold of $N+1$ coefficients, not the FFT, dominates the cost.

\paragraph{In the deep tail} $\gamma\to1$ and the tilt flattens the per-category laws. For the
probability-mass statistic the effect is exact: at $\gamma=1$ the factor $(1-\gamma)$ in the proof
of Proposition~\ref{prop:logconcave} vanishes, $\pi_k(i)\propto e^{\nu i}$, and every category carries
the \emph{same} truncated geometric law, whatever $p_k$ is, with mean $N/K$. A geometric law has
standard deviation comparable to its mean, so $V\approx K(N/K)^2$ and $\sqrt V\approx N/\sqrt K$. For LLR,
Stirling gives $e^{\ell_k(i)+\nu i}\propto e^{(1+\nu')i}/\sqrt{2\pi i}$ at $\gamma=1$ --- geometric up
to the $\sqrt i$ factor --- and the same picture holds. Measured at $N=10^4$ with uniform $x_0$:
$\sqrt V=2471$, $2362$, $1989$, $1558$, $1060$ for $K=3,5,10,20,50$ (LLR), against
$N/\sqrt K=5774$, $4472$, $3162$, $2236$, $1414$; the ratio $0.43$--$0.75$ reflects $\gamma<1$ and the
truncation at $N$, which bites for small $K$.

So in the regime this codebase exists for, $Q$ grows linearly in $N$ and the advantage over the tree
is a constant factor rather than an asymptotic one; it is largest for many categories, where
$\sqrt V$ shrinks like $1/\sqrt K$ while $M$ does not change.

\section{Why the approximate de-Poissonization is not enough}
\label{sec:approx}

The gamma search (\texttt{gamma\_search.h}) uses the same Poissonized weights but never undoes the
Poissonization: it minimizes $\log(\rho^{-N}\prod_kZ_k)+\log\sinck$, which omits even the
normalization $P(T=N)\approx1/\sqrt{2\pi V}$. That is harmless there, because $\gamma$ only
conditions the series --- the tilt cancels exactly --- so an $O(1/V)$ error in the objective merely
moves $\gamma$ slightly.

For the terms, $M_S(z_n)$ \emph{is} the answer, and the natural approximation --- the Gaussian
(saddlepoint) evaluation of the circle integral,
$\coef N F\approx F(\rho)\rho^{-N}/\sqrt{2\pi V}$, i.e.\ keeping only the neighbourhood of $q=0$ ---
fails twice. At $n=0$ it carries the usual $O(1/V)$ relative error of a saddlepoint approximation
\citep{daniels1954}, straight into $p$. At larger $n$ it fails outright: the phases
$e^{-in\theta_k(i)}$ decorrelate the categories, the integrand is no longer concentrated near
$q=0$, and the Gaussian picture has nothing to approximate. Measured against the exact terms:

\begin{center}
\small
\begin{tabular}{lcccc}
\toprule
instance & $\gamma$ & $\sqrt V$ & relative error at $n=0$ & worst relative error, $n<512$ \\
\midrule
LLR $N=100$, $K=20$, $x_0$ at the mode      & 0.137 & 10.8  & $9.7\times10^{-4}$ & $0.45$ \\
LLR $N=1000$, $K=7$, uniform $x_0$          & 0.996 & 399   & $3.7\times10^{-2}$ & $0.46$ \\
PMF $N=1000$, $K=5$, $x_0$ at the mode      & 0.287 & 37.4  & $1.2\times10^{-4}$ & $1.00$ \\
\bottomrule
\end{tabular}
\end{center}

The exact version costs the $Q$-point sum instead of one point, and removes both errors.

\section{An aside: selecting $\gamma$ by the same Poissonisation}
\label{sec:gamma}

The gamma search just referred to (\texttt{multinomial/gamma\_search.h}) is worth deriving in full,
since it is the other place in this codebase where dropping the multinomial constraint turns an
expensive combinatorial object into a cheap independent one -- and, unlike the terms, it gets away
with never restoring the constraint at all.

\texttt{fill\_gamma}'s job is to choose the real tilt $\gamma$ that \texttt{fill\_cells} bakes into
every reward before the series runs. As \texttt{error\_bounds.tex} establishes, the natural
objective is the log-MGF of the \emph{smoothed} statistic,
$\Psi(\gamma):=\log M_{S+U}(\gamma)=\log M_S(\gamma)+\log\sinck(\pi\gamma/T)$, and its stationary
point is the saddlepoint condition $E_\gamma[S+U]=0$. The grid search that used to solve this
(still available as \texttt{fill\_gamma\_grid}) evaluated $\Psi$ exactly, via the same $K$-fold
constrained convolution as the terms, at 128 points refined over several rounds. Poissonising
$\Psi$ turns each evaluation into the one-pass computation of \S\ref{sec:poisson}: with
$w_k(i;\gamma,\nu)=e^{\ell_k(i)+\nu i}$ now carrying an explicit $\gamma$ inside $\ell_k(i)$ as well
(recall $\ell_k(i)=\log p_k(i)-\gamma a_k(i)$ once \texttt{fill\_cells} has run; here we treat
$\gamma$ as the free variable and undo that baking-in for the derivation) and $Z_k(\gamma,\nu)=\sum_iw_k(i;\gamma,\nu)$,
define
\[
  A(\gamma,\nu):=\sum_k\log Z_k(\gamma,\nu),\qquad
  \Psi_{\mathrm{Pois}}(\gamma):=\min_\nu\big[A(\gamma,\nu)-\nu N\big]+\log N!+\log\sinck(\pi\gamma/T).
\]
$A(\gamma,\cdot)-\nu N$ is convex in $\nu$ (a cumulant generating function minus a linear term), so
its minimiser $\nu(\gamma)$ is exactly the constraint $\sum_kE_k[i]=N$ of \S\ref{sec:poisson} --
the same tilt \texttt{gamma\_moments} solves for by Newton bisection, warm-started from the
previous $\gamma$'s solution.

\begin{proposition}[$\Psi_{\mathrm{Pois}}$ is convex, and its derivatives are free]
\label{prop:gammaderiv}
Write $x=\pi\gamma/T$, $g=\log\sinck$, and let $E_k[\cdot]$, $\Var_k[\cdot]$, $\Cov_k(\cdot,\cdot)$
denote moments under $w_k(\cdot;\gamma,\nu(\gamma))/Z_k$. Then
\[
  \Psi_{\mathrm{Pois}}'(\gamma)=-\sum_kE_k[a_k]+\tfrac{\pi}{T}g'(x),
  \qquad
  \Psi_{\mathrm{Pois}}''(\gamma)=\sum_k\Var_k[a_k]-\frac{\big(\sum_k\Cov_k(a_k,i)\big)^2}{\sum_k\Var_k[i]}+\Big(\tfrac{\pi}{T}\Big)^2g''(x)\;\ge\;0.
\]
\end{proposition}
\begin{proof}
Write $V(\gamma):=\min_\nu[A(\gamma,\nu)-\nu N]=A(\gamma,\nu(\gamma))-\nu(\gamma)N$. Since
$\nu(\gamma)$ is an unconstrained interior minimiser, the envelope theorem gives
$V'(\gamma)=\partial_\gamma A(\gamma,\nu(\gamma))=\sum_k\partial_\gamma\log Z_k=-\sum_kE_k[a_k]$
(the $\nu'(\gamma)$ term drops out because $\partial_\nu[A-\nu N]=0$ at $\nu(\gamma)$), which gives
the first identity together with the sinc contribution.

For the second, differentiate the two first-order quantities
$f_1(\gamma,\nu):=\sum_kE_k[i]-N$ and $f_2(\gamma,\nu):=-\sum_kE_k[a_k]$
along the implicit curve $\nu=\nu(\gamma)$ defined by $f_1=0$. Direct computation from
$w_k=\exp(-\gamma a_k(i)+\log p_k(i)+\nu i)$ gives, for each $k$,
$\partial_\gamma E_k[i]=-\Cov_k(a_k,i)$, $\partial_\nu E_k[i]=\Var_k[i]$,
$\partial_\gamma E_k[a_k]=-\Var_k[a_k]$, $\partial_\nu E_k[a_k]=\Cov_k(a_k,i)$ (each a standard
exponential-tilting identity: differentiating $\log Z_k$ once more brings down another centred
factor). The implicit function theorem applied to $f_1=0$ gives
$\nu'(\gamma)=-\partial_\gamma f_1/\partial_\nu f_1=\sum_k\Cov_k(a_k,i)/\sum_k\Var_k[i]$, and then
\[
  \Psi_{\mathrm{Pois}}''(\gamma)-\Big(\tfrac{\pi}{T}\Big)^2g''(x)
  =\frac{d}{d\gamma}f_2(\gamma,\nu(\gamma))
  =\partial_\gamma f_2+\partial_\nu f_2\cdot\nu'(\gamma)
  =\sum_k\Var_k[a_k]-\frac{\big(\sum_k\Cov_k(a_k,i)\big)^2}{\sum_k\Var_k[i]}.
\]
Writing $R=\sum_ka_k(I_k)$ and $I=\sum_kI_k$ for the (independent, under the Poissonised measure)
category totals, this right-hand side is exactly $\Var(R)-\Cov(R,I)^2/\Var(I)$, the residual
variance $\min_\beta\Var(R-\beta I)$ of the best linear predictor of $R$ from $I$ -- non-negative
because a variance is. $\sinck$'s own log is convex on $\mathbb R$ (it is $-\log$ of a normalised
completely monotone kernel), so $g''\ge0$ too, and $\Psi_{\mathrm{Pois}}''\ge0$ follows.
\end{proof}

This is a compressed restatement, in the Poissonised setting, of the general fact underlying every
saddlepoint method: cumulant generating functions are convex, and their derivatives are moments.
What makes it useful here is that both moments in
Proposition~\ref{prop:gammaderiv} come out of \emph{the same single pass} over the $K(N+1)$ cells
that \texttt{gamma\_moments} already runs to solve for $\nu(\gamma)$ -- accumulating $\sum_i xw_k(i)$,
$\sum_ix^2w_k(i)$ for $x\in\{i,a_k(i)\}$ and their cross term costs one extra multiply-add per cell
over the moment $\nu$ already needed, so the derivative is free in the sense that matters: no
second evaluation of the objective, no finite-difference step size to choose.

\paragraph{The search.} \texttt{gamma\_newton\_search} is a safeguarded Newton iteration (Newton
where it behaves, bisection where it does not, in the style of Numerical Recipes' \texttt{rtsafe})
on $\Psi_{\mathrm{Pois}}'$ over the caller's bracket, unchanged at $[-2,2]$. A subtlety: the step
size cannot be the stopping test here, because $\Psi_{\mathrm{Pois}}''$ spans several orders of
magnitude across the bracket -- the sinc term alone contributes $(\pi/T)^2g''(0)=(\pi/T)^2/12$ near
$\gamma=0$ but only $\sim1/\gamma^2$ once $\pi\gamma/T$ exceeds one -- so plain Newton can take
steps far smaller than the remaining distance to the root without indicating it should stop. Only
the width of the bracket is used to decide that, exactly as \texttt{eps\_gamma} already specifies.
Once the bracket has shrunk below \texttt{eps\_gamma}, one further safeguarded Newton refinement
of the last iterate is applied (clipped to that final bracket): it is accurate to the square of the
bracket's own error, essentially for free, and matters because the series length is acutely
sensitive to exactly where in a narrow bracket $\gamma$ lands.

\paragraph{Why the geometric shortcut used for the terms does not apply here.} The frequency loop
of \S\ref{sec:impl} advances $e^{-in\theta_k(i)}$ by a constant complex factor because that phase
has no data-dependent rescaling in it. Every sum inside \texttt{gamma\_moments} (and, before it,
the FFT-tree's own \texttt{find\_nu\_gamma\_vec}) is instead a log-sum-exp whose numerically
stabilising pivot (\texttt{max\_exp}) is chosen per $\gamma$ -- and per SIMD lane, when several
$\gamma$ are evaluated together. The exponents actually formed are $w_k(i)/e^{\text{pivot}(\gamma)}$
rather than $w_k(i)$ itself, and that data-dependent shift destroys the geometric relation between
the exponentials at neighbouring $\gamma$. The fix is the one above -- evaluate $\gamma$ fewer
times -- not to make each evaluation itself incremental.

\paragraph{Measured.} On the instances of Table~\ref{tab:verify} and a few others, at
\texttt{eps\_rel}$=10^{-6}$:

\begin{center}
\small
\begin{tabular}{rrlrrrrr}
\toprule
$N$ & $K$ & stat & $\gamma$, grid $\to$ Newton & evaluations & time, grid $\to$ Newton (ms) & speedup & $|\Delta p|/p$ \\
\midrule
1000 & 7  & LLR & $0.997229\to0.996335$ & $128\to6$  & $53.1\to4.0$   & $13\times$ & $6.2\times10^{-9}$ \\
100  & 20 & LLR & $0.993197\to0.993762$ & $128\to6$  & $16.9\to1.2$   & $15\times$ & $4.0\times10^{-12}$ \\
100  & 20 & LLR & $0.147644\to0.137001$ & $128\to8$  & $34.6\to1.4$   & $25\times$ & $3.3\times10^{-7}$ \\
1000 & 5  & LLR & $0.985135\to0.982263$ & $128\to7$  & $66.1\to3.1$   & $21\times$ & $9.7\times10^{-6}$ \\
1000 & 5  & PMF & $0.334089\to0.286481$ & $128\to7$  & $42.3\to2.5$   & $17\times$ & $6.9\times10^{-5}$ \\
500  & 50 & LLR & $0.831948\to0.827904$ & $128\to6$  & $182.5\to13.3$ & $14\times$ & $3.4\times10^{-14}$ \\
2000 & 5  & LLR & $0.993197\to0.991124$ & $128\to7$  & $83.0\to6.0$   & $14\times$ & $6.3\times10^{-6}$ \\
\bottomrule
\end{tabular}
\end{center}
Both searches disagree on $\gamma$ itself by $5\times10^{-4}$--$5\times10^{-2}$ (the surrogate's
saddlepoint is exact only to $O(1/N)$), yet the resulting p-values agree to $10^{-14}$--$10^{-5}$:
as \S\ref{sec:approx} already notes, $\gamma$ only conditions the series -- the tilt cancels out of
$p$ exactly -- so this level of disagreement is expected and harmless. What is not free is the time:
$13$--$25\times$ per instance here, and $12\times$ aggregated over a sweep of 747 instances spanning
$N\in[20,1000]$, $K\in[2,20]$ and both statistics in earlier measurement.

\section{The evaluator}
\label{sec:impl}

The implementation is \texttt{src/multinomial/fourier\_poisson.h}, selected by
\texttt{options->poisson = 1}, which overrides \texttt{fft\_precision} and
\texttt{use\_fft\_precompute} for the series terms and never builds the merge-factor plan.

\paragraph{Plan} (\texttt{poisson\_plan\_t}, built once per \texttt{series()}): the tilt $\nu$ and
variance $V$; the normalized weights $\pi_k(i)$ for all $K(N+1)$ cells; $Q$ with its certificate; the
bit-reversed destination of each count $i\bmod Q$; the $Q/2$ roots; the extractor
$\omega^{qN}/Q$ in the code's sign convention; and the magnitude
$\texttt{resN\_exp}=\exp(\log N!+\sum_k\log Z_k-\nu N)$ in quad precision.

\paragraph{Per frequency vector.} As in the other evaluators, one pass carries $\mathtt{STRIDE}=8$
frequencies in SIMD lanes. For each category: zero $2Q$ lanes of bins; fold
$\pi_k(i)e^{-in\theta_k(i)}$ into the bit-reversed bin of $i\bmod Q$, with the phase produced by the
recurrence of \texttt{fourier.h} (per-thread \texttt{phase} times shared \texttt{lane}), so no
transcendental function is evaluated in the loop; run the iterative FFT, whose input is already
permuted and whose first level uses root one only; and multiply the result into the running product.
Finally one dot product with the extractor gives the unit of~\eqref{eq:main}, multiplied by
$\texttt{resN\_exp}$ and $\sinck(z_n)$ as in every other evaluator. Every lane shares the fold
positions, the roots and the product structure, so every operation is a whole-vector one.

\paragraph{Cost and memory.} Per frequency the work is
\[
  K\Big[(N+1)\;+\;\tfrac{Q}{2}\log_2Q\;+\;Q\Big]\;+\;Q
\]
elementary complex operations, against $K\cdot\mathrm{calls}(M)$ merges~\eqref{eq:calls} for the
tree, and there is a single forward transform per category instead of a descent and an ascent.
The plan takes $8K(N+1)+O(N+Q)$ bytes and each thread $4Q$ vectors ($256Q$ bytes) of scratch, against
$16K\cdot\mathrm{calls}(M)$ bytes and about $3M$ vectors per thread for the tree.

\paragraph{Threads.} The evaluator is a drop-in thread function: thread $t$ takes the frequencies
$n\equiv t\pmod{\texttt{threads}}$, so that the phase recurrence has a constant step and survives
chunk boundaries. The plan is read-only and shared. The deviations of Table~\ref{tab:verify} are
the same to all printed digits at 8, 3 and 1 threads.

\paragraph{A persistent pool.} \texttt{series()} calls whichever evaluator is selected once per
chunk of \texttt{threads * vecs\_per\_thread * STRIDE} frequencies -- typically some tens to a few
hundred times per run, occasionally into the thousands for a slowly converging instance. The
pre-existing dispatcher, \texttt{eval\_fourier\_parallel}, called
\texttt{pthread\_create} and \texttt{pthread\_join} \texttt{threads} times on \emph{every} call,
which is unnecessary work repeated every chunk: nothing about a worker's job changes chunk to
chunk except which frequencies it takes next. \texttt{fourier\_pool\_t} instead creates the
worker threads once per \texttt{series()} call and reuses them for every chunk, through a
textbook reusable barrier -- a generation counter and two condition variables, one workers wait on
for the next chunk and one the caller waits on for completion -- so only the fields that actually
change (\texttt{start}, \texttt{end}, \texttt{warm}) are touched between chunks; everything else
(\texttt{mult}, \texttt{res\_arr}, the plan, the phase tables, even the evaluator function pointer)
is written into each worker's persistent \texttt{arguments\_t} once, at pool creation.

In isolation (no evaluator work, just the synchronisation), waking 8 parked threads costs about
$20\,\mu$s against $61\,\mu$s to create and join 8 fresh ones -- a $3\times$ reduction in pure
overhead. Against the real evaluators, that overhead is usually a small fraction of one chunk's
actual work (a few percent at the instance sizes in Table~\ref{tab:thru}, lost in run-to-run
noise), but it dominates at the smallest instances, where a chunk's work is itself only tens of
microseconds:
\begin{center}
\small
\begin{tabular}{rrlrrrr}
\toprule
$N$ & $K$ & evaluator & threads & create/join ($\mu$s/chunk) & pool ($\mu$s/chunk) & speedup \\
\midrule
20  & 2 & Poisson & 8 & 50.3 & 25.5 & $2.0\times$ \\
20  & 3 & Poisson & 8 & 56.5 & 29.5 & $1.9\times$ \\
50  & 3 & Poisson & 8 & 67.7 & 40.3 & $1.7\times$ \\
20  & 2 & tree     & 8 & 64.9 & 40.3 & $1.6\times$ \\
\bottomrule
\end{tabular}
\end{center}
The pool also removes the two \texttt{malloc}/\texttt{free} pairs (\texttt{arguments\_t} and the
\texttt{pthread\_t} array) that \texttt{eval\_fourier\_parallel} repeated on every call; that cost
is not visible in the timings above but disappears with the same change. At \texttt{threads}$=1$
the pool skips its synchronisation machinery entirely and calls the evaluator directly, matching
the previous sequential path.

\section{Verification}
\label{sec:verify}

Table~\ref{tab:verify} compares the new evaluator with the FFT tree on instances chosen to span both
statistics, small and large $K$, the deep tail and moderate p-values, and $Q$ from 128 to 8192. Term
by term the two agree to within the aliasing tolerance and rounding; the maximum deviation is the
same to all printed digits at 8, 3 and 1 threads. At a fixed term budget, with stopping and
extrapolation disabled, the p-values agree to between $10^{-15}$ and $2.3\times10^{-12}$ relative.
The one row where the term deviation reaches $7.5\times10^{-11}$ is the one whose
$Q$ was accepted at a certificate of $7.4\times10^{-11}$, exactly as Lemma~\ref{lem:dominate}
predicts, and its p-value still agrees to $2.3\times10^{-12}$.

\begin{table}[t]
\centering
\small
\begin{tabular}{rrlcrrrccc}
\toprule
$N$ & $K$ & stat & $x_0$ & $\gamma$ & $\sqrt V$ & $Q$ & $D(Q)$ & $\max_n|\Delta t_n|/t_0$ & $|\Delta p|/p$ \\
\midrule
100  & 20 & LLR & mode           & 0.137 & 10.8  & 128  & $8.7\times10^{-15}$ & $1.6\times10^{-14}$ & $1.1\times10^{-14}$ \\
100  & 20 & LLR & unif, $k^5$    & 0.994 & 35.1  & 512  & $6.1\times10^{-15}$ & $9.4\times10^{-14}$ & $9.5\times10^{-14}$ \\
1000 & 7  & LLR & unif, $k^2$    & 0.996 & 399.2 & 4096 & $9.7\times10^{-14}$ & $2.5\times10^{-13}$ & $2.6\times10^{-13}$ \\
1000 & 7  & PMF & unif, $k^2$    & 0.997 & 327.9 & 4096 & $3.1\times10^{-14}$ & $1.5\times10^{-12}$ & $1.5\times10^{-12}$ \\
1000 & 5  & PMF & mode           & 0.287 & 37.4  & 512  & $1.9\times10^{-13}$ & $9.4\times10^{-14}$ & $4.3\times10^{-14}$ \\
1000 & 3  & LLR & unif           & 0.984 & 247.4 & 2048 & $1.1\times10^{-14}$ & $3.9\times10^{-13}$ & $3.7\times10^{-13}$ \\
1000 & 2  & LLR & unif           & 0.985 & 232.6 & 1024 & exact & $1.9\times10^{-13}$ & $1.2\times10^{-13}$ \\
500  & 50 & LLR & unif           & 0.828 & 52.6  & 512  & $1.9\times10^{-14}$ & $2.1\times10^{-13}$ & $1.9\times10^{-13}$ \\
50   & 3  & LLR & unif, $k^{1/2}$& 0.475 & 9.8   & 128  & exact & $7.7\times10^{-15}$ & $1.0\times10^{-15}$ \\
5000 & 10 & LLR & unif           & 0.995 & 994.1 & 8192 & $7.4\times10^{-11}$ & $7.5\times10^{-11}$ & $2.3\times10^{-12}$ \\
\bottomrule
\end{tabular}
\caption{Poissonized evaluator against the FFT tree. Term deviations over the first 512 frequencies
(256 at $K=50$, 128 at $N=5000$), identical at 8, 3 and 1 threads; p-values at the same fixed term
budget with \texttt{eps\_rel = 0} and \texttt{extrapolate = 0}. ``unif'' is $x_0=N/K$, ``mode'' is
$x_0\approx N\mathbf p$, and $k^s$ gives $p_k\propto k^s$ ($s=1$ when omitted). ``exact'' marks
$Q>\max(N,(K-1)N)$, where no alias exists and no certificate is needed.}
\label{tab:verify}
\end{table}

\section{The statistic's interval, without the $O(KN^2)$ recursion}
\label{sec:interval}

Before the series can run, \texttt{fill\_interval} must locate the range $[s_{\min},s_{\max}]$ of
$S(x)-S(x_0)=\sum_kg_k(x_k)-g_k(x_{0,k})$ over integer $x\ge0$, $\sum_kx_k=N$, where $g_k$ is the
category-$k$ contribution to whichever statistic is in use (the LLR, a general power-divergence, or
the pmf log-density). It did this with a $K(N{+}1)$-cell $\min$/$\max$-plus convolution, the
combinatorial analogue of the tree evaluator's own recursion -- $O(KN^2)$, and, at $N=10^4$,
$K=50$, the single largest item in setup (\S\ref{sec:speed} below). Both endpoints admit a
much cheaper closed form.

\begin{proposition}[$g_k$ is concave]
\label{prop:gconcave}
For every real $\lambda$ (the power-divergence family, LLR at $\lambda=0$), $g_k(x)=
-\tfrac1{\lambda(\lambda+1)}x\big((x/Np_k)^\lambda-1\big)$ is strictly concave on $x>0$:
\[
  g_k''(x)=-(Np_k)^{-\lambda}x^{\lambda-1}<0.
\]
The LLR limit $\lambda\to0$, $g_k(x)=-x\log(x/Np_k)$, has $g_k''(x)=-1/x<0$ directly. For the pmf
statistic, $g_k(x)=x\log p_k-\log\Gamma(x{+}1)$ is concave because $\log\Gamma$ is convex on
$(0,\infty)$ --- the defining property, by the Bohr--Mollerup theorem, that makes $\Gamma$ the
unique log-convex extension of the factorial to the reals --- so $g_k$, an affine function of $x$
minus a convex one, is concave.
\end{proposition}
\begin{proof}
Differentiate the power-divergence case twice; the factor $\lambda(\lambda+1)$ introduced by the
prefactor cancels exactly against the one produced by differentiating $x^{\lambda+1}$ twice, so the
sign is uniformly negative with no special case at $\lambda\in(-1,0)$. The pmf case needs no
derivative at all: log-convexity of $\Gamma$ on $(0,\infty)$ is exactly the Bohr--Mollerup
characterization of it.
\end{proof}

\begin{proposition}[exact minimum, $O(K)$]
\label{prop:intervalmin}
$S(x)=\sum_kg_k(x_k)$, being concave and separable, attains its minimum over the polytope
$\{x\ge0,\sum x_k=N\}$ exactly at one of its $K$ vertices $Ne_k$ (``all $N$ observations in category
$k$''): $s_{\min}=\min_k\,g_k(N)-r_0$, where $r_0=S(x_0)$.
\end{proposition}
\begin{proof}
Writing any $x$ in the polytope as the convex combination $x=\sum_k(x_k/N)(Ne_k)$ of the vertices
and applying concavity, $S(x)\ge\sum_k(x_k/N)S(Ne_k)\ge\min_kS(Ne_k)$; the vertices are themselves
feasible, so the minimum over the whole polytope equals the minimum over just the $K$ vertices.
$S(Ne_k)=g_k(N)+\sum_{j\ne k}g_j(0)=g_k(N)$ since $g_j(0)=0$. Shifting by $x_0$, the vertex value
of $S(x)-S(x_0)$ is $g_k(N)-r_0$, and both $g_k(N)$ and $r_0$ are already available from
\texttt{mult->matrix} and \texttt{mult->r0} -- no new evaluation of the reward function is needed.
\end{proof}

\begin{proposition}[maximum, power-divergence family: closed-form bound, $O(1)$]
\label{prop:intervalmax}
For the power-divergence family (any real $\lambda$, LLR included), $s_{\max}\le-r_0$, with
equality of the \emph{continuous} relaxation (dropping integrality on $x$).
\end{proposition}
\begin{proof}
Substituting $x_k=Np_k$ into $g_k'$ gives $g_k'(Np_k)=-1/(\lambda+1)$ for every $k$ -- the same value
regardless of $p_k$ -- so $x=Np$ satisfies the stationarity condition of the equality-constrained
Lagrangian with a single shared multiplier, and by concavity (Proposition~\ref{prop:gconcave}) that
stationary point is the global maximum of the continuous relaxation. Since the base of the power
ratio is $1$ there, $g_k(Np_k)=0$ for every $k$, so the continuous maximum of $S$ is exactly $0$ and
of $S(x)-S(x_0)$ is exactly $-r_0$. The true, integer-constrained maximum can only be smaller, since
the integer lattice is a subset of the relaxed domain. (This is the familiar fact that a
power-divergence statistic is non-negative and vanishes only when the observation matches its
expectation -- no search is needed to find it.)
\end{proof}

\paragraph{Maximum, pmf statistic: GreedyModeFind, $O(K\log K)$.} Proposition~\ref{prop:intervalmax}
does not apply to the pmf statistic: $g_k'(x)=\log p_k-\psi(x{+}1)$ has no closed-form inverse, and
$x=Np$ is not in general a stationary point of $\sum_kg_k$ (the pmf is not itself a divergence
normalised to vanish at the mean, only asymptotically so). But $S(x)=\log P(x)-\log N!$, so
maximizing $S$ over the integer polytope is exactly the problem of finding a \emph{mode} of
$\mathrm{Multinomial}(N,p)$ -- and \citet{whitehendy2010}'s \texttt{GreedyModeFind} solves that
exactly in $O(K\log K)$. It starts from the unique mode $w_k=\lfloor ap_k\rfloor$ of a nearby,
trivially computed distribution (their Theorem 1, for $a=N+K/2$), then walks frequency by frequency
towards $N$; at each step, incrementing (or decrementing) whichever category $k$ maximises
$p_k/(x_k{+}1)$ (resp.\ $x_k/p_k$) strictly increases $P(x)$ and preserves optimality among reachable
vectors (their Theorems 2--3). A binary max-heap keyed on that per-category score turns ``pick the
best category, update its score'' into an $O(\log K)$ operation, for $O(K\log K)$ total (we use only
their single-mode algorithm, \texttt{GreedyModeFind1}; the tie-tracking needed to enumerate the full
set of joint modes is unnecessary here, since only the value $S(x^\ast)$ at one mode $x^\ast$ is
wanted). $s_{\max}=S(x^\ast)-r_0$ is then exact, not merely a bound.

\paragraph{Measured.} Both formulas were checked against \texttt{fill\_interval}'s DP directly (not
just through the series) across uniform, skewed ($p_k\propto k^2$), and deliberately dominant-category
($p_1\gg p_{k>1}$) distributions, $N$ up to $500$, $K$ up to $50$, and $\lambda\in\{-0.5,-0.3,0,2/3,1,2\}$
plus the pmf statistic: $s_{\min}$ agrees with the DP to floating-point noise ($<10^{-8}$ relative) in
every case, including the $\lambda\ne0$ family's exact-in-theory formula and the pmf statistic's
GreedyModeFind, and $s_{\max}$ for the power-divergence family is always a valid upper bound (never
below the DP's exact value), with slack small relative to $|s_{\min}|$ -- consistent with $s_{\min}$
being the binding constraint on $W=\max(s_{\max},-s_{\min})$ in every case checked, including the
adversarial dominant-category distributions. Isolated from the rest of setup, at $K=100$: the DP's
time scales as expected with $N$ (about $4\times$ per doubling, from $0.10$\,s at $N=2000$ to $9.1$\,s
at $N=20{,}000$ for the pmf statistic), while GreedyModeFind finishes in about $2\,\mu$s regardless of
$N$ -- roughly five orders of magnitude faster at $N=20{,}000$. \S\ref{sec:speed} below reports the
effect on the whole pipeline.

\section{Measured speedup}
\label{sec:speed}

All timings are on an Apple M5 (4 performance and 6 efficiency cores) with \texttt{gcc-15 -Ofast
-march=native}, $\mathtt{STRIDE}=8$ lanes through SIMDe's AVX-512 emulation on NEON, the best of three
runs, and uniform $x_0=N/K$ with $p_k\propto k$ --- the shape of the benchmark suite, and the deep
tail --- unless stated.

\subsection{Per frequency}

Table~\ref{tab:thru} isolates the evaluators: the cost of one frequency, excluding plan construction,
at one thread and at the default eight.

\begin{table}[tp]
\centering
\small
\begin{tabular}{rrlrrrrrrrr}
\toprule
 & & & & & \multicolumn{3}{c}{1 thread ($\mu$s)} & \multicolumn{3}{c}{8 threads ($\mu$s)} \\
\cmidrule(lr){6-8}\cmidrule(lr){9-11}
$N$ & $K$ & stat & $\sqrt V$ & $Q$ / $M$ & tree & Poisson & speedup & tree & Poisson & speedup \\
\midrule
100  & 3  & LLR & 25.4  & 256 / 256   & 9.7    & 3.2   & $3.0\times$ & 2.9    & 1.1   & $2.7\times$ \\
100  & 5  & LLR & 24.4  & 256 / 256   & 16.5   & 5.2   & $3.2\times$ & 4.2    & 1.7   & $2.6\times$ \\
100  & 10 & LLR & 20.2  & 256 / 256   & 32.8   & 9.9   & $3.3\times$ & 7.0    & 2.9   & $2.4\times$ \\
100  & 20 & LLR & 15.7  & 128 / 256   & 66.3   & 9.6   & $6.9\times$ & 13.6   & 2.8   & $4.8\times$ \\
100  & 50 & LLR & 11.1  & 128 / 256   & 165.4  & 23.7  & $7.0\times$ & 33.4   & 5.0   & $6.7\times$ \\
100  & 3  & PMF & 26.1  & 256 / 256   & 10.0   & 3.3   & $3.1\times$ & 2.9    & 1.1   & $2.7\times$ \\
100  & 5  & PMF & 25.1  & 256 / 256   & 16.6   & 5.3   & $3.2\times$ & 4.3    & 1.7   & $2.6\times$ \\
100  & 10 & PMF & 20.9  & 256 / 256   & 33.2   & 10.2  & $3.3\times$ & 7.1    & 2.9   & $2.4\times$ \\
100  & 20 & PMF & 16.8  & 256 / 256   & 66.5   & 20.0  & $3.3\times$ & 13.6   & 4.8   & $2.9\times$ \\
100  & 50 & PMF & 12.9  & 128 / 256   & 166.0  & 23.9  & $6.9\times$ & 34.0   & 5.0   & $6.8\times$ \\
\midrule
1000 & 3  & LLR & 247.4 & 2048 / 2048 & 108.8  & 39.4  & $2.8\times$ & 22.5   & 7.8   & $2.9\times$ \\
1000 & 5  & LLR & 236.5 & 2048 / 2048 & 181.6  & 63.9  & $2.8\times$ & 36.9   & 12.3  & $3.0\times$ \\
1000 & 10 & LLR & 198.8 & 2048 / 2048 & 361.6  & 124.5 & $2.9\times$ & 72.9   & 23.6  & $3.1\times$ \\
1000 & 20 & LLR & 155.5 & 2048 / 2048 & 724.0  & 247.1 & $2.9\times$ & 145.8  & 46.1  & $3.2\times$ \\
1000 & 50 & LLR & 105.5 & 1024 / 2048 & 1804.7 & 286.1 & $6.3\times$ & 359.2  & 55.4  & $6.5\times$ \\
1000 & 3  & PMF & 250.4 & 2048 / 2048 & 108.8  & 39.2  & $2.8\times$ & 22.4   & 7.9   & $2.8\times$ \\
1000 & 5  & PMF & 237.0 & 2048 / 2048 & 181.1  & 63.6  & $2.8\times$ & 36.9   & 12.3  & $3.0\times$ \\
1000 & 10 & PMF & 195.0 & 2048 / 2048 & 361.4  & 124.3 & $2.9\times$ & 73.5   & 23.9  & $3.1\times$ \\
1000 & 20 & PMF & 150.4 & 2048 / 2048 & 720.4  & 246.8 & $2.9\times$ & 145.2  & 46.2  & $3.1\times$ \\
1000 & 50 & PMF & 102.1 & 1024 / 2048 & 1808.2 & 286.2 & $6.3\times$ & 354.8  & 55.4  & $6.4\times$ \\
\midrule
$10^4$ & 3  & LLR & 2471 & 32768 / 32768 & 2251.8  & 775.6  & $2.9\times$  & 816.4   & 209.0  & $3.9\times$ \\
$10^4$ & 5  & LLR & 2362 & 32768 / 32768 & 3805.0  & 1289.3 & $3.0\times$  & 1314.8  & 335.8  & $3.9\times$ \\
$10^4$ & 10 & LLR & 1989 & 16384 / 32768 & 7636.2  & 1303.7 & $5.9\times$  & 2600.2  & 277.1  & $9.4\times$ \\
$10^4$ & 20 & LLR & 1558 & 16384 / 32768 & 15288.7 & 2597.5 & $5.9\times$  & 5242.5  & 553.4  & $9.5\times$ \\
$10^4$ & 50 & LLR & 1060 & 8192 / 32768  & 38406.6 & 3451.0 & $11.1\times$ & 13277.9 & 668.0  & $19.9\times$ \\
$10^4$ & 3  & PMF & 2495 & 32768 / 32768 & 2298.4  & 803.3  & $2.9\times$  & 800.7   & 225.6  & $3.6\times$ \\
$10^4$ & 5  & PMF & 2358 & 32768 / 32768 & 3821.7  & 1320.1 & $2.9\times$  & 1339.6  & 353.8  & $3.8\times$ \\
$10^4$ & 10 & PMF & 1937 & 32768 / 32768 & 7648.2  & 2559.5 & $3.0\times$  & 2673.8  & 664.6  & $4.0\times$ \\
$10^4$ & 20 & PMF & 1488 & 16384 / 32768 & 15319.4 & 2567.9 & $6.0\times$  & 5317.0  & 571.8  & $9.3\times$ \\
$10^4$ & 50 & PMF & 997  & 8192 / 32768  & 38235.0 & 3452.7 & $11.1\times$ & 13513.3 & 683.7  & $19.8\times$ \\
\midrule
\multicolumn{11}{l}{\emph{moderate p-values}, $x_0$ at the mode} \\
100  & 20 & LLR & 10.8  & 128 / 256   & 66.4   & 9.7   & $6.9\times$  & 13.7   & 2.8   & $4.8\times$ \\
1000 & 5  & PMF & 37.4  & 512 / 2048  & 180.7  & 15.1  & $12.0\times$ & 36.9   & 4.0   & $9.3\times$ \\
1000 & 20 & LLR & 33.2  & 256 / 2048  & 721.0  & 34.5  & $20.9\times$ & 146.0  & 7.3   & $20.1\times$ \\
\bottomrule
\end{tabular}
\caption{Cost of one series term, FFT tree (precompute plan) against the Poissonized evaluator,
both dispatched through the persistent thread pool of \S\ref{sec:impl} -- the same path
\texttt{series()} actually uses.}
\label{tab:thru}
\end{table}

Three patterns stand out.
\begin{itemize}
  \item \textbf{Deep tail with $Q=M$: about $3\times$} ($2.4$--$4.0\times$). With the same transform
        size the new evaluator is still faster, because a butterfly is cheaper than a merge and there
        is one pass per category instead of a descent and an ascent.
  \item \textbf{Each halving of $Q$ roughly doubles the gain}: $4$--$9\times$ at $Q=M/2$ and, at
        $Q=M/4$ ($N=10^4$, $K=50$), $11\times$ at one thread widening to $20\times$ at eight. Because
        $Q$ moves in powers of two, similar instances can land a factor of two apart: LLR and PMF at
        $N=10^4$, $K=10$ have $\sqrt V$ within $3\%$ of each other, but the certificate accepts
        $Q=16384$ for LLR and not for PMF, whose geometric per-category laws have the heavier tails.
  \item \textbf{Away from the tail}, $Q\ll M$ and the fold dominates: $7$--$21\times$ at one thread,
        $5$--$20\times$ at eight.
\end{itemize}
At $N=100$ the eight-thread ratios are still somewhat compressed relative to one thread (e.g.\
$2.7\times$ against $3.0\times$ for LLR $K=3$) because at $1$--$4\,\mu$s per frequency, even the
persistent pool's synchronisation cost (\S\ref{sec:impl}, about $2.5\,\mu$s per chunk at 8 threads)
is not entirely negligible next to the work; both evaluators pay it, and it is far smaller than the
$\sim$$8\,\mu$s per chunk the old per-chunk \texttt{pthread\_create}/\texttt{join} cost here.

At $N=10^4$, $K=50$ the pattern reverses: the tree/Poisson speedup is \emph{larger}, not smaller,
at eight threads than at one ($19.9\times$ against $11.1\times$ for LLR, $19.8\times$ against
$11.1\times$ for PMF) -- the widest thread-count gap in the table, and the opposite of the
small-instance compression just above. The likely reason is the per-thread footprint gap of
\S\ref{sec:memplan}: $12.5$\,MB of scratch per thread for the tree at $N=10^4$ against $2$--$8$\,MB
for the Poissonized evaluator, so eight tree threads compete for memory bandwidth by the time eight
Poissonized ones do not, and the tree's parallel scaling suffers exactly where its plan is largest.

\subsection{Memory and plan construction}
\label{sec:memplan}

\begin{center}
\small
\begin{tabular}{rrrrrr}
\toprule
 & & \multicolumn{2}{c}{plan (MB)} & \multicolumn{2}{c}{plan build (ms)} \\
\cmidrule(lr){3-4}\cmidrule(lr){5-6}
$N$ & $K$ & tree & Poisson & tree & Poisson \\
\midrule
1000   & 50 & 27.3  & 0.41 & 5.3   & 1.4  \\
$10^4$ & 3  & 35.2  & 1.06 & 6.9   & 3.4  \\
$10^4$ & 10 & 117.5 & 1.21 & 21.2  & 7.0  \\
$10^4$ & 20 & 235.0 & 1.98 & 41.8  & 12.5 \\
$10^4$ & 50 & 587.5 & 4.08 & 115.3 & 15.0 \\
\bottomrule
\end{tabular}
\end{center}
Per-thread scratch is $12.5$\,MB for the tree at $N=10^4$ against $256Q$ bytes, $2$--$8$\,MB, for the
Poissonized evaluator. The Poisson plan build includes the certificate's scalar evaluations.

\subsection{End to end}

Finally the whole pipeline --- setup, gamma search, plan and adaptive series, now through the
persistent thread pool of \S\ref{sec:impl}, the robust stopping rule of \S\ref{sec:stopping} and the
$O(K)$/$O(K\log K)$ statistic interval of \S\ref{sec:interval} in place of \texttt{fill\_interval}'s
$O(KN^2)$ recursion --- with the default options (\texttt{extrapolate = 2}), \texttt{max\_iter =
100000} and eight threads, at \texttt{eps\_rel} of $10^{-3}$ and $10^{-6}$: 72 runs over
$N\in\{100,1000,5000,10^4\}$. Across all of them the series time \emph{per term} improves by a
factor between $2.0$ and $22.1$, with median $3.0$. The stopping rule of \S\ref{sec:stopping} makes
the tree and Poisson evaluators report the \emph{same} term count in 71 of the 72 runs (the one
exception differs by $6$ terms out of $670$). A selection, with speedups of the series alone and of
the whole run:

\begin{center}
\small
\begin{tabular}{rrlcrrrrc}
\toprule
$N$ & $K$ & stat & \texttt{eps\_rel} & setup (s) & series, tree $\to$ Poisson (s) & series & total & terms \\
\midrule
1000   & 5  & LLR & $10^{-3}$ & 0.0024 & $0.0353\to0.0126$  & $2.8\times$  & $2.5\times$  & 888 \\
1000   & 50 & LLR & $10^{-6}$ & 0.0221 & $0.1478\to0.0266$  & $5.6\times$  & $3.4\times$  & 332 \\
5000   & 5  & LLR & $10^{-6}$ & 0.0200 & $7.913\to2.364$    & $3.4\times$  & $3.3\times$  & 14915 \\
5000   & 10 & LLR & $10^{-6}$ & 0.0354 & $0.818\to0.130$    & $6.3\times$  & $5.0\times$  & 783 \\
5000   & 50 & PMF & $10^{-6}$ & 0.177  & $2.165\to0.142$    & $15.3\times$ & $6.4\times$  & 327 \\
$10^4$ & 5  & LLR & $10^{-3}$ & 0.0384 & $0.834\to0.276$    & $3.0\times$  & $2.7\times$  & 546 \\
$10^4$ & 20 & LLR & $10^{-3}$ & 0.183  & $1.274\to0.164$    & $7.8\times$  & $3.8\times$  & 173 \\
$10^4$ & 50 & LLR & $10^{-6}$ & 0.308  & $7.298\to0.330$    & $22.1\times$ & $10.3\times$ & 333 \\
\midrule
1000   & 5  & PMF, mode & $10^{-6}$ & 0.0030 & $0.540\to0.060$  & $9.0\times$  & $8.6\times$  & 12025 \\
1000   & 20 & LLR, mode & $10^{-3}$ & 0.0096 & $0.203\to0.013$  & $15.7\times$ & $9.1\times$  & 1164 \\
\bottomrule
\end{tabular}
\end{center}

Term counts are given once per row since the two evaluators now agree exactly in every row shown.
Setup is now dominated by \texttt{fill\_rewards} ($O(KN)$: filling the $K(N{+}1)$-cell reward
matrix) and the gamma search, rather than by locating the statistic's range: replacing
\texttt{fill\_interval}'s $O(KN^2)$ recursion with the closed-form bound and GreedyModeFind of
\S\ref{sec:interval} cuts setup at $N=10^4,K=50$ from $1.45$\,s to $0.31$\,s, and at
$N=5000,K=50$ (PMF) from $0.75$\,s to $0.18$\,s. With setup no longer the larger cost even at
$N=10^4$, the end-to-end gain across the table now tracks the series speedup much more closely:
$1.7$--$14.6\times$ over the full 72-run sweep, against $2.1$--$4.4\times$ when \texttt{fill\_interval}
was still the setup path.

Some instances at $\texttt{eps\_rel}=10^{-6}$ genuinely do not converge within the $100{,}000$-term
budget under either stopping criterion -- LLR and PMF at $K=3$, $x_0=N/K$, are the clearest cases in
this sweep. This is not a regression from either fix: re-running the plain term-ratio rule alone
(\texttt{extrapolate}$=0$) on the same instances with the cap raised to $200{,}000$ terms shows it
does not converge there either: the slow convergence is a property of these instances, not of
either stopping criterion.

\section{Robust stopping for \texttt{extrapolate} $\ge2$}
\label{sec:stopping}

\texttt{extrapolate}$\,\ge2$ stops the series on either of two independent conditions, whichever is
satisfied first: the pre-existing term-ratio rule (\texttt{CONV\_COUNT} consecutive terms with
$|{\rm term}_n/p_{n-1}|<\texttt{eps\_rel}/\texttt{CONV\_COUNT}$), or a check on a second,
shallow extrapolation of the same partial sums.

That second extrapolation runs the same DQELG algorithm as the reporting one (\S\ref{sec:impl}),
but with a much shallower table (depth 8 rather than 50 -- deep columns are exactly where
catastrophic cancellation lives), and its raw output is smoothed by a running median of the last 5
values. Convergence is a question about that smoothed sequence: has it stayed within
\texttt{eps\_rel} of itself over a trailing window of $\max(\texttt{CONV\_COUNT},\,i/2)$ terms. A
single outlier term cannot flip this test, because by the time it would matter it has aged out of
both the median and the window. When this check is the one that stops the series, \texttt{pval}
and \texttt{err\_rel\_ext} are also drawn from this same smoothed statistic and its trailing-window
spread, rather than from the deep reporting table -- the value reported is always the one that
justified stopping.

\paragraph{Stability under perturbation.} Repeating every run in this section with every term's
magnitude perturbed by $-1$, $0$ and $+1$ ulp of \texttt{lgac[N]} (a relative change of
$10^{-13}$--$10^{-12}$), at \texttt{eps\_rel}$=10^{-6}$:

\begin{center}
\small
\begin{tabular}{llrr}
\toprule
instance & evaluator (threads) & terms ($-1$/0/$+1$ ulp) & $|\Delta p|/p$ ($-1$,$+1$ vs.\ 0) \\
\midrule
LLR $N=1000$, $K=3$ & tree (8)    & 200000 / 200000 / 200000 & $9\times10^{-13}$, $9\times10^{-13}$ \\
PMF $N=1000$, $K=5$ & tree (8)    & 15341 / 15341 / 15341    & $9\times10^{-13}$, $8\times10^{-12}$ \\
LLR $N=100$, $K=5$  & tree (8)    & 106249 / 106249 / 106249 & $6\times10^{-14}$, $6\times10^{-14}$ \\
PMF $N=100$, $K=3$  & tree (8)    & 200000 / 200000 / 200000 & $6\times10^{-14}$, $6\times10^{-14}$ \\
PMF $N=1000$, $K=5$ & Poisson (8) & 15341 / 15341 / 15341    & $5\times10^{-10}$, $6\times10^{-9}$ \\
PMF $N=1000$, $K=5$ & Poisson (3) & 15341 / 15341 / 15341    & $3\times10^{-11}$, $2\times10^{-10}$ \\
PMF $N=1000$, $K=5$ & Poisson (1) & 15341 / 15341 / 15341    & $7\times10^{-12}$, $9\times10^{-13}$ \\
LLR $N=1000$, $K=3$ & Poisson (1) & 200000 / 200000 / 200000 & $9\times10^{-13}$, $9\times10^{-13}$ \\
\bottomrule
\end{tabular}
\end{center}

Every instance gives the same term count for all three perturbations, at every evaluator and thread
count tried, and the reported $p$ agrees to $10^{-9}$--$10^{-13}$ relative -- the size of the
perturbation itself, not an amplified reaction to it.

\paragraph{The window check firing.} On LLR, $N=1000$, $K=7$, $p_k\propto k^2$,
$\texttt{eps\_rel}=10^{-3}$, the window check stops the series at $92$ terms, well before the
term-ratio rule would (which needs $224$). The result is still trustworthy: $4.4\times10^{-5}$
relative error against a $65536$-term reference, within the requested $10^{-3}$, and
\texttt{err\_rel\_ext}$=8.6\times10^{-4}$ correctly reports a wider margin than the $224$-term run's
$2.0\times10^{-6}$ would. On LLR, $N=1000$, $K=20$, $x_0$ at the mode, $\texttt{eps\_rel}=10^{-3}$,
the same $1164$-term stopping point is reproduced under all three perturbations, and the reported
$p$ agrees to $9\times10^{-13}$ across them, honestly labelled by a larger
\texttt{err\_rel\_ext} of $3$--$8\times10^{-2}$ in this instance -- a case where
$\texttt{eps\_rel}=10^{-3}$ is optimistic for the number of terms actually available, and the
diagnostic says so rather than reporting a falsely precise-looking number.

\section{Summary}
\label{sec:summary}

\begin{itemize}
  \item Dropping the multinomial constraint makes the $K$ categories independent; a $Q$-point
        trapezoid rule for Cauchy's integral on the saddlepoint circle restores it exactly up to
        aliasing (Proposition~\ref{prop:cauchy}).
  \item The aliasing of every term is bounded by the aliasing at frequency zero
        (Lemma~\ref{lem:dominate}), and that is measured exactly by comparing $Q$ with $2Q$
        (Proposition~\ref{prop:cert}). The measurement is a certificate for the probability-mass
        statistic (Proposition~\ref{prop:logconcave}) and matches the true aliasing to two or three digits in
        every measured case, including LLR.
  \item $Q$ scales with the spread of the tilted total count: about $10\sqrt N$ away from the tail,
        and about $N/\sqrt K$ in the deep tail, where the tilt makes the per-category laws
        geometric.
  \item Against the FFT tree: $p$ agrees to $10^{-15}$--$2.3\times10^{-12}$; a single term is about
        $3\times$ cheaper in the deep tail when $Q=M$, $4$--$20\times$ when $Q$ is $M/2$ or $M/4$, and
        $4$--$22\times$ away from the tail; the plan is $30$--$140\times$ smaller.
  \item The evaluator is opt-in (\texttt{options->poisson}); it never builds the tree's plan.
  \item The same construction selects $\gamma$: the Poissonized log-MGF is convex, its derivatives
        come from the envelope theorem out of the moment pass already run for the tilt $\nu$
        (Proposition~\ref{prop:gammaderiv}), and a safeguarded Newton iteration on it is
        $13$--$25\times$ faster than the 128-point grid it replaces, per instance measured.
  \item A persistent thread pool replaces per-chunk thread creation in \texttt{series()}: up to
        $2\times$ at the smallest instances, where synchronisation overhead is otherwise comparable
        to the work itself; negligible where the per-chunk work already dominates.
  \item The \texttt{extrapolate}$\,\ge2$ stopping rule combines the term-ratio rule with a check on
        a shallow, median-smoothed extrapolation's trailing-window stability, and reports from that
        same robust statistic whenever it is the one that stops the series; verified stable to
        $10^{-9}$--$10^{-13}$ under a one-ulp perturbation of every term, across both evaluators and
        every thread count tried, and confirmed to stop a genuinely ringing instance early
        ($224\to92$ terms) without loss of accuracy (\S\ref{sec:stopping}).
  \item Locating the statistic's range, previously an $O(KN^2)$ dynamic program, is now $O(K)$ for
        the power-divergence family (concavity gives an exact vertex-based minimum and a
        closed-form maximum bound) and $O(K\log K)$ for the probability-mass statistic (an exact
        maximum via \citeauthor{whitehendy2010}'s multinomial-mode algorithm), verified against the
        DP across skewed and dominant-category distributions and cutting setup time by
        $4$--$5\times$ at $N=10^4$ (\S\ref{sec:interval}).
\end{itemize}

\begin{thebibliography}{9}

\bibitem[Bornemann(2011)]{bornemann2011}
F.~Bornemann.
\newblock Accuracy and stability of computing high-order derivatives of analytic functions by
Cauchy integrals.
\newblock \emph{Foundations of Computational Mathematics}, 11(1):1--63, 2011.

\bibitem[Daniels(1954)]{daniels1954}
H.~E.~Daniels.
\newblock Saddlepoint approximations in statistics.
\newblock \emph{Annals of Mathematical Statistics}, 25(4):631--650, 1954.

\bibitem[Hoggar(1974)]{hoggar1974}
S.~G.~Hoggar.
\newblock Chromatic polynomials and logarithmic concavity.
\newblock \emph{Journal of Combinatorial Theory, Series B}, 16(3):248--254, 1974.

\bibitem[Jacquet and Szpankowski(1998)]{jacquetszpankowski1998}
P.~Jacquet and W.~Szpankowski.
\newblock Analytical depoissonization and its applications.
\newblock \emph{Theoretical Computer Science}, 201(1--2):1--62, 1998.

\bibitem[Levin(1981)]{levin1981}
B.~Levin.
\newblock A representation for multinomial cumulative distribution functions.
\newblock \emph{The Annals of Statistics}, 9(5):1123--1126, 1981.

\bibitem[Lyness and Moler(1967)]{lynessmoler1967}
J.~N.~Lyness and C.~B.~Moler.
\newblock Numerical differentiation of analytic functions.
\newblock \emph{SIAM Journal on Numerical Analysis}, 4(2):202--210, 1967.

\bibitem[Trefethen and Weideman(2014)]{trefethenweideman2014}
L.~N.~Trefethen and J.~A.~C.~Weideman.
\newblock The exponentially convergent trapezoidal rule.
\newblock \emph{SIAM Review}, 56(3):385--458, 2014.

\bibitem[White and Hendy(2010)]{whitehendy2010}
W.~T.~J.~White and M.~D.~Hendy.
\newblock A fast and simple algorithm for finding the modes of a multinomial distribution.
\newblock \emph{Statistics and Probability Letters}, 80(1--2):63--68, 2010.

\end{thebibliography}

\end{document}
